\newpage
\section{Wymagania funkcjonalne}

System sterowania dronem przy użyciu interfejsu BCI wymaga spełnienia szeregu funkcji, które zapewnią prawidłowe działanie wszystkich komponentów, od rejestracji sygnałów EEG po sterowanie dronem w czasie rzeczywistym. Poniżej przedstawiono szczegółowe wymagania dotyczące infrastruktury przetwarzania danych, która stanowi podstawe całego systemu i odpowiada za przesyłanie, przetwarzanie oraz magazynowanie informacji w spójny i bezpieczny sposób.

\subsection{Infrastruktura przetwarzania danych}

\subsubsection{Fizyczna infrastruktura klastra}
\begin{itemize}
    \item Klaster powinien składać się z pięciu fizycznych komputerów połączonych w sieć.
    \item Maszyny powinny umożliwiać równoległe przetwarzanie danych EEG.
    \item Klaster musi być dostępny zdalnie poprzez VPN.
    \item Na każdym komputerze powinien działać system Ubuntu Server.
    \item Klaster powinien być skalowalny, umożliwiając potencjalne dodanie nowych węzłów w przyszłości.
\end{itemize}
\subsubsection{Przesyłanie napływających danych EEG i treningowych}
\begin{itemize}
    \item System powinien umożliwiać przesyłanie strumieniowe danych EEG i danych treningowych pomiędzy headsetem, klastrem i modułem uczenia maszynowego.
    \item Dane powinny docierać w czasie rzeczywistym, bez utraty próbek i z zachowaniem integralności sygnału.
    \item System powinien umożliwiać kolejkowanie danych i obsługę wielu źródeł jednocześnie.
    \item System powinien umożliwiać wstępną weryfikację jakości napływających danych i oznaczanie podejrzanych próbek.
\end{itemize}
\subsubsection{Przetwarzanie danych w strumieniu czasu rzeczywistego}
\begin{itemize}
    \item System powinien filtrować i oczyszczać dane EEG w czasie rzeczywistym, eliminując artefakty i szumy.
    \item Powinna być możliwa ekstrakcja cech dla algorytmów uczenia maszynowego po odebraniu danych.
    \item System powinien wspierać wysyłanie wyników przetwarzania do dalszych komponentów (np. modeli ML, dashboardu) w czasie rzeczywistym.
\end{itemize}
\subsubsection{Magazynowanie danych surowych i przetworzonych}
\begin{itemize}
    \item System musi przechowywać surowe dane EEG w formacie umożliwiającym późniejszą analizę i retraining modeli.
    \item Przetworzone dane i wektory cech powinny być magazynowane w sposób uporządkowany 
    \item Dane powinny być przechowywane na wielu węzłach klastra, co zapewnia redundancję i bezpieczeństwo w przypadku awarii pojedynczej maszyny.
\end{itemize}
\subsubsection{Orkiestracja procesów przetwarzania i treningu modeli}
\begin{itemize}
    \item System powinien automatycznie zarządzać przepływem danych między komponentami.
    \item Powinien umożliwiać harmonogramowanie treningu modeli uczenia maszynowego w określonych interwałach lub po zgromadzeniu nowych danych treningowych.
\end{itemize}
\subsubsection{Monitoring i wizualizacja stanu systemu}
\begin{itemize}
    \item System powinien umożliwiać podgląd jakości napływających danych EEG i statusu przetwarzania w czasie rzeczywistym.
    \item Powinien monitorować obciążenie klastra, wykorzystanie pamięci i CPU poszczególnych maszyn.
    \item System powinien wizualizować wyniki przetwarzania, błędy i statystyki wydajności.
    \item Powinna istnieć możliwość generowania raportów i alarmów w przypadku wykrycia nieprawidłowości w działaniu systemu.
\end{itemize}
\subsection{Preprocessing i ekstrakcja cech ?????}
\subsection{Machine Learning}
\subsection{Headset EEG}
\subsection{Dron}
\subsection{Aplikacja}